\setAuthor{Erkki Tempel}
\setRound{piirkonnavoor}
\setYear{2019}
\setNumber{G 5}
\setDifficulty{5}
\setTopic{Magnetism}

\prob{Osake magnetväljas}
Osake laenguga $q$, massiga $m$ ja kiirusega $v$ siseneb magnetvälja induktsiooniga $B$. Kui lai peab minimaalselt olema magnetvälja ala $l$, et osake liiguks pärast magnetväljast väljumist esialgsele liikumissuunale vastassuunas?
\begin{center}
  \includegraphics[width = 0.4\linewidth]{2019-v2g-05-yl.pdf}
\end{center}\hint

\solu
Kui osake siseneb magnetvälja, siis mõjub osakesele Lorenzi jõud $F_L = qvB\sin{\alpha}$. Osake siseneb magnetvälja risti magnetväljaga ning hakkab Lorentzi jõi tõttu liikuma mööda ringjoont.

Ringjoone raadius $r$ on minimaalne siis, kui osake siseneb magnetvälja alt ning väljub ülevalt servast. Mööda ringjoont liikudes mõjub osakesele tsentrifugaaljõud
\[ F_T = \frac{mv^2}{r}\]
Lorentzi jõud ja tsentrifugaaljõud on võrdsed , seega saame avaldada ringjoone raadiuse $r$
\[ F_L = F_T \quad\quad\Rightarrow\quad\quad qvB = \frac{mv^2}{r} \quad\quad\Rightarrow\quad\quad r = \frac{mv^2}{qvB} \]
Seega on magnetvälja ala minimaalne laius $l$
\[ l = 2r = \frac{2mv^2}{qvB} \]

 \vspace{-20pt}
  \begin{center}
    \includegraphics[width=0.5\textwidth]{2019-v2g-05-sol.pdf}
  \end{center}
  \vspace{-20pt}\probend